\section{Encoding arbitrary common graphs}
\label{gen:section}

This section gives an exact reduction to the complete-sector extension
of Synchronized Planarity, which is solved in Section~\ref{cse:section}
(Theorem~\ref{cse:main}). In particular, the
compilation below does not assume tractability of arbitrary simultaneous
ordering, arbitrary triple-orientation constraints, or arbitrary finite-domain
constraint satisfaction.

\begin{theorem}[General face-list compilation]
\label{gen:compilation}
Let $I=(G_1,G_2)$ have connected colours and biconnected union $U$.  Suppose
that every P-node $\mu$ of every nontrivial common block is supplied with a
cyclic order $\pi_\mu$ such that every simultaneous embedding induces either
$\pi_\mu$ or its reversal at that node.  There is a deterministic
polynomial-time construction of an instance of the complete-sector
Synchronized Planarity language which is feasible if and only if $I$ admits
a \SEFE.  Both witness transformations are polynomial.

The construction uses only complete-sector observations, unary domain
restrictions, partial bijections, and affine equations on Boolean variables,
in addition to the ordinary graph, pipes, and Q-cells.  With
$N=|V|+|E_1|+|E_2|$, its output has size $O((N+1)^3)$, and a direct implementation of the construction takes
$O((N+1)^4)$ time.
\end{theorem}

The reduction theorem uses the
SPQR embedding parameterization, the source-labelled embedding-tree
construction, Lemma~\ref{rep:lem:two-partitions}, and the repair
Theorem~\ref{rep:thm:repair}.

\subsection{Normalization and the fixed common P orders}
\label{gen:porders}

The normalization of Section~\ref{sec:normalization} first makes the colours
connected and splits at union cutvertices.  Let $\mu$ be a P-node of a
nontrivial common block $K$, with poles $p,q$.  Its directions are the
$\{p,q\}$-bridges of $K$ (Section~\ref{gen:sectors}); a direction is
\emph{non-real} unless it is the edge $pq$, and its \emph{interior} is its
vertex set minus $\{p,q\}$.  Two non-real directions are \emph{linked} if
$U$ contains a path from an interior vertex of one to an interior vertex of
the other whose internal vertices avoid $V(K)$; equivalently, some
component of $U-V(K)$ is adjacent to both interiors, or a private edge
joins them.  The \emph{union-link graph} of $\mu$ has the non-real
directions as nodes and the linked pairs as edges.  This is the union-link
graph of~\cite{bkr}, restricted to non-real directions.  Linked directions
are consecutive in the P order of every \SEFE{}~\cite[Lemma~2]{bkr}.

The common-cycle separation-pair decomposition of
Bl\"asius--Karrer--Rutter~\cite[Lemma~4 and Theorem~2]{bkr} then produces
an equivalent family of total linear size in which the union-link graph of
every common P-node is connected.  Here the split components at a
separation pair $\{u,v\}$ of $U$ are the components of $U-\{u,v\}$ with
their edges to $u,v$; an existing or added edge $uv$ is not counted as a
split component carrying a common path.  Without this convention a factor
together with its added edge $uv$ would again separate a common cycle, and
the decomposition would not terminate.  The qualification ``non-real'' is
necessary: a real common pole edge has no interior and would be isolated in
a link graph that included it.  By~\cite[Lemma~3]{bkr}, in a feasible
instance every connected union-link graph is a spanning path or cycle.  A
spanning path fixes the order of the non-real directions up to reversal.
A real common pole edge, if present, must be inserted in the closing gap
of the path, because every other gap is forced by a link that would have
to cross it~\cite[Lemma~2]{bkr}; neither~\cite{bkr} nor~\cite{fpr} states
this placement explicitly.  A spanning cycle fixes the order when there is
no real common pole edge and otherwise causes rejection.  These rules
supply the reference cyclic order $\pi_\mu$.  Any other connected link
graph is a valid rejection.  See also~\cite[Lemma~4(c)]{fpr}.
Throughout, $\pi_\mu$ is the cyclic order at the pole $p$; at $q$ the same
embedding induces the reversed order.

For clarity about intermediate edge identities, a fresh common two-edge
path can replace a newly added common pole-to-pole marker.  Restricting a
\SEFE{} to each factor and contracting a common path in another factor gives
the forward witness map.  Conversely, the common markers permit the factor
drawings to be glued in disjoint pole-to-pole regions, after which the
markers are deleted.  Each colour remains connected: every colour component
of a factor meets at least one pole, and the marker connects the poles.
Each nontrivial augmented union factor is biconnected.  Subdivision makes
parallel marker edges harmless if a simple-graph representation is desired.
Perform it after the decomposition and retain the already computed
reference orders.  The orders $\pi_\mu$ are computed \emph{before} this
subdivision, with every marker treated as a real common pole edge in the
rules above; after subdivision a marker would be an unlinked non-real
direction, and the rules would wrongly reject.  Subdivided markers are not
subjected to a new link connectivity test.
The stronger BKR preprocessing that removes certain further two-pole split
components is not needed in this section.

\subsection{Attachment sets, bridges, and reference sectors}
\label{gen:sectors}

For a nontrivial common block $K$, let $B_i(K)$ be its unique containing
colour-$i$ block.  For $Z\in\cc(U-V(K))$ and a colour lobe
$Y\in\cc(B_i(K)-V(K))$, write
\[
 A(K,Z)=N_U(Z)\cap V(K),\qquad
 A_i(K,Y)=N_{B_i(K)}(Y)\cap V(K).
\]
Each lobe lies in a unique $Z$.  Biconnectivity gives at least two distinct
attachments for every $Z$ and every lobe.  Common bridges have only one
face and need no additional face rows.  They and isolated common vertices
remain present in the ordinary source-tree construction and in the final
repair.

A $\{p,q\}$-bridge of a biconnected graph is a component after deleting
$p,q$, with its incident pole edges, or the real edge $pq$ if present.
There is a canonical P-node at $p,q$ exactly when there are at least three
such bridges.  For a common P-node $\mu$ with poles $p,q$, write
\[
 \pi_\mu=(e_0,e_1,\ldots,e_{k-1}),\qquad k\geq 3.
\]
Name its reference sector $j$ by the consecutive pair
$(e_j,e_{j+1})$, with indices modulo $k$.  The Boolean $r_\mu$ specifies
whether the common P order is the reference order or its reversal.
Reflection changes the oriented corner used by a named sector, not its
name.  Common R-sources also receive global reflection variables relative
to fixed reference embeddings.

\paragraph{The fixed sign convention for represented tree copies.}
Give each represented tree copy a bit $\varepsilon$: it is $0$ for an
original red tree or an auxiliary $X$ tree, and $1$ for an original blue
tree or its comparison $K'$ tree.  If $\rho_B$ is its original block
frontier, the rotation in the ordinary constraint graph is
\[
 \widehat\rho=\rev^{\varepsilon}(\rho_B).
\]
For the same three labelled leaves, the median identity therefore reads
\begin{equation}
 \orcyc_{\rho_B}(a,b,c)
   =\orcyc_{\widehat\rho}(a,b,c)\oplus\varepsilon
   =\orcyc_{\widehat\rho_m}(a_m,b_m,c_m)\oplus\varepsilon.
 \label{gen:copy-sign}
\end{equation}
Here $m$ is their median and $a_m,b_m,c_m$ are its three directions.
All side observations use the original directed cycle's outgoing and
incoming labelled leaves; they do not acquire a second sign from the
direction in which a doubled marked route happens to be traversed.

For an R-copy, choose its represented reference rotation to be
$\rev^{\varepsilon}$ of its original source reference.  Its Q-cell flip
then equals the original native reflection bit.  For a P-copy, do the
same with its marked reference list.  If a list
$(a_0,\ldots,a_{h-1})$ is reversed while retaining $a_0$ as its first
entry, represented sector $j$ corresponds to original sector
$-j-1\pmod h$.  Thus the sector relabeling is the identity for
$\varepsilon=0$ and $j\mapsto-j-1\pmod h$ for $\varepsilon=1$.
This also swaps the two sector names when $h=2$.  The fixed maps in the
proof below include this explicit relabeling and the chosen pole
convention.  Equation~\eqref{gen:copy-sign}, rather than an unspecified
blue offset, is used whenever an original triple bit is compiled.

\begin{lemma}[Faces at a P-pair]
\label{gen:p-faces}
In every embedding of $K$ with P order $\pi_\mu$ up to reversal, its
faces incident with both $p$ and $q$ are in bijection with the $k$ reference
sectors of $\mu$.
\end{lemma}

\begin{proof}
A nonedge bridge has connected interior and at most one face incident with
both poles.  Indeed, two pole-to-pole arcs in distinct such faces would
form a Jordan curve through only the poles with bridge edges on both
sides, disconnecting the interior after deletion of $p,q$.  In the
embedding of a bond of bridges, every other bridge occupies that unique
pole-to-pole exterior region.  Gluing the bridges in their cyclic order
therefore creates exactly one common face between every consecutive pair.
The edge bridge is handled as one direction of the bond.  Every face
incident with both poles arises in this way.  Internal changes to a bridge
may alter the boundary path of such a face, but cannot alter its named
sector.
\end{proof}

\subsection{Common directions at a colour source}

Let $B$ be a biconnected supergraph of $K$ (a colour block or an auxiliary
graph), let $\mu$ be a common P-node of $K$ with poles $p,q$, and let $J$
be a $\{p,q\}$-bridge of $B$; when $B$ has a P-source at $p,q$, these are
its expansions.  Two common $\mu$-directions contained in $J$ are
\emph{colour-linked} if $B$ contains a path between interior vertices of
the two directions that lies in $J-\{p,q\}$ and has no internal vertex in
$V(K)$; equivalently, some component of $B-V(K)$ inside $J$ is adjacent to
both interiors, or a private edge of $B$ joins them.  The
\emph{colour-link graph} of $J$ has the common $\mu$-directions in $J$ as
nodes and the colour-linked pairs as edges.

\begin{lemma}[Projection to common bridge groups]
\label{gen:groups}
Let $K$ be a biconnected subgraph of a biconnected graph $B$, and let $\nu$
be a P-source of $B$ with poles $p,q$.  If at most one pole belongs to $K$,
all edges of $K$ lie in one source expansion.  If at least three source
expansions meet $K$, then $p,q$ are poles of a common P-node $\mu$.
The source expansions meeting $K$ partition the common $\mu$-directions
into groups, each of which is a cyclic interval in every extending common
P order.  Their cyclic order is the projection of that common order.
\end{lemma}

\begin{proof}
Deleting at most one vertex from $K$ leaves it connected, proving the first
statement.  If both poles belong to $K$, each component of
$K-\{p,q\}$ lies in one source expansion; a real edge $pq$ is also
indivisible.  Conversely, if an expansion contains a vertex of $K$ other
than its poles, connectivity of $K-p$ connects that vertex to $q$ inside
the expansion, and connectivity of $K-q$ connects it to $p$.  Thus every
expansion meeting $K$ contains one or more whole common bridges.
Three such expansions imply at least three common bridges and the stated
common P-node.

The interior of one colour bridge is connected.  Paths in it between
common directions, broken at successive visits to $K$, give a connected
colour-link graph on those directions.  Each link forces its two common
directions to be consecutive, because the link lies in one common face.
Equivalently, all other colour bridges lie in the unique pole-to-pole
exterior face of this bridge.  Its common directions therefore form one
cyclic interval, and the order of the colour bridges is the projected
order of these intervals.
\end{proof}

If a proposed reference order has a noninterval group, it cannot extend to
a planar embedding of $B$ and is rejected.  Otherwise the projected group
order supplies the prescribed complete marked order at every source P-copy
with at least three marked directions.  This prescribes only its marked
subset; every unmarked direction can still be placed in any permitted gap
and in any order within that gap.

Two different common blocks cannot both display a cycle at the same
P-source: both displayed cycles would contain both poles, whereas distinct
common blocks share at most one vertex.

\subsection{An injective read of an entire two-pole face}

\begin{lemma}[Complete face read]
\label{gen:complete-read}
Suppose that $A_i(K,Y)=\{p,q\}$ and that $p,q$ are the poles of a
common P-node $\mu$.  The face of $K$ containing $Y$ can be represented
exactly by a variable $t_Y$ with domain $\{0,\ldots,k-1\}$ using one of
the following forms: a constant; an injective relabeling of a Boolean
variable; or an injective relabeling of one complete source-sector
category.  The form and all maps are computable in polynomial time.
\end{lemma}

\begin{proof}
The lobe $Y$ is an entire component of $B_i(K)-\{p,q\}$.  A path from
it to a vertex of $K-\{p,q\}$ would first create an additional attachment
to $K$, contrary to the hypothesis.  Thus $Y$ is a private pole-to-pole
bridge with nonempty interior.

Partition the common $\mu$-directions according to their containing
colour pole-to-pole bridges, and let $h$ be the number of nonempty groups.
Suppose first that $h\geq 2$.  Together with $Y$ there are at least
$h+1\geq 3$ bridges, so the colour graph has a P-source $\nu$ at these
same poles.  Its marked directions are precisely the $h$ groups meeting
$K$, and its direction for $Y$ is unmarked.  Each complete source sector
maps to the common gap between the last direction of one interval group
and the first direction of the next.  The $h$ resulting gaps are distinct.
Call this injection $\iota$.

For $h\geq 3$, the source frame is synchronized with $r_\mu$, with the
fixed reference and pole convention offset, by the chain rows of
Section~\ref{gen:sync} (see the paragraph after~\eqref{gen:chain-rows}).
Consequently $t_Y$ is the fixed injective relabeling $\iota$ of its framed
source category, independently of reflection choices.

For $h=2$, the two anchors alone imply no frame equality.  Let
$(a_0,a_1)$ be the ordered reference pair of the source copy
(Section~\ref{gen:marking}), let $G_0$ be the group whose colour bridge
contains $a_0$, and let $G_1$ be the other group.  Let $\iota(0)$ be the
reference sector between the last direction of $G_0$ and the first
direction of $G_1$ in $\pi_\mu$, and let $\iota(1)$ be the other sector
between the two groups.  Let $c$ be the two-sector category of the
dart of $Y$ at that copy and $f$ the frame of that copy.  Introduce a Boolean $b$
with
\begin{equation}
 b=c\oplus f\oplus r_\mu\oplus\delta,
 \qquad t_Y=\iota(b),
 \qquad \delta=\varepsilon\oplus\sigma,
 \label{gen:two-gauge}
\end{equation}
where $\varepsilon$ is the copy bit of~\eqref{gen:copy-sign}, and
$\sigma=0$ if the read is made at the pole $p$ at which $\pi_\mu$ is read
and $\sigma=1$ if it is made at $q$.  To verify the constant, take
$\varepsilon=0$ and a read at $p$.  By Definition~\ref{cse:definition},
$c\oplus f=0$ exactly when the dart of $Y$ lies after $a_0$ and before
$a_1$ in the rotation at $p$.  Going forward from the bridge of $G_0$ to
that of $G_1$ crosses the common gap $\iota(0)$ if the common order at $p$
is $\pi_\mu$ ($r_\mu=0$), and $\iota(1)$ if it is reversed.  A read at $q$
sees the reversed common order, and a represented copy with
$\varepsilon=1$ sees the reversed rotation with the same frame; each adds
$1$.  The first relation in~\eqref{gen:two-gauge} is affine and the
second is a partial bijection.  Formula~\eqref{gen:two-gauge} retains the
independent two-anchor gauge explicitly.

Now let $h=1$.  A real common edge $pq$ cannot occur in this case:
it would itself be a separate colour bridge.  Thus every common direction
has an interior.  All common directions lie in one colour bridge.  Its
colour-link graph is connected: a path in the bridge minus the poles,
broken at successive visits to $K$, gives a link-graph path between any
two common directions.  In a planar embedding it is a spanning path or
cycle of the common cyclic order.  A cycle is impossible here.  The
private bridge $Y$ would occupy some common pole-to-pole face.  A colour
link across that sector is a crosscut of the face disk joining interiors
of its two boundary pole-to-pole paths, and hence separates $p$ from $q$
inside the disk.  The required colour pole-to-pole path in $Y$ could not
cross it.
Thus the link graph is a spanning path.  The gap between its endpoints is
the only common sector available to a private pole-to-pole bridge.  Pin
$t_Y$ to that sector.  Its reference name is determined by the graph and
$\pi_\mu$, independently of the embedding.  If the colour-link graph is not
a spanning path whose linked pairs are consecutive in $\pi_\mu$, then no
planar embedding of $B_i(K)$ induces $\pi_\mu$ or its reversal, and the
construction rejects.  This is correct, because every \SEFE{} induces
$\pi_\mu$ or its reversal.

These expressions record the actual face in every planar block embedding.
They impose no other choice of order or position inside that face.  Real
private chords $pq$ can be additional unmarked bridges, but they are not
lobes and do not affect the argument.  The same read may be made at the
other pole $q$ with the same injection $\iota$: for $h=2$ this is the
case $\sigma=1$ of~\eqref{gen:two-gauge}; for $h\geq3$ the pole
convention is part of the fixed reference offset of the chain rows; for
$h=1$ the face is a constant.  No further correction is applied.  Thus the
statement applies to either pole copy at which a later basis-side
observation occurs.
\end{proof}

\subsection{Synchronizing all common R and P choices}
\label{gen:sync}

For every occurrence of a biconnected supergraph $B$ of $K$, whether an
actual colour block or an auxiliary graph, synchronize its induced
common embedding with the global source variables of $K$.

For a common R-node $\mu$, select three actual common darts at a skeleton
vertex behind three distinct skeleton directions.  Their median in
$\Theta_B(v)$ is an R-source.  A P median with other pole $q$ would
separate those three darts after deleting $v,q$.  This contradicts the
3-connectivity of the common R-skeleton: $q$ affects at most one incident
expansion, or one other skeleton vertex, and the other two selected
directions remain connected through the R-skeleton after those deletions.
Thus $r_\mu$ equals a native R flip plus a computable reference constant.
Add that affine row.

For a common P-node $\mu$, write its reference order as
$\pi_\mu=(e_0,\ldots,e_{k-1})$, and select one real common dart
$a_j$ behind each direction $e_j$ at one pole.  A linear family of
triple-orientation equations suffices to synchronize the full order.

\begin{lemma}[Chain characterization of a cyclic order]
\label{gen:chain-lemma}
Let $\alpha=(a_0,\ldots,a_{k-1})$ be a reference cyclic order and
$r\in\Ftwo$.  A cyclic order $\rho$ equals $\alpha$ when $r=0$ and
its reversal when $r=1$ if and only if
\begin{equation}\label{gen:order-chain}
 \orcyc_\rho(a_0,a_i,a_{i+1})
  =\orcyc_\alpha(a_0,a_i,a_{i+1})\oplus r
 \quad (1\le i\le k-2).
\end{equation}
\end{lemma}
\begin{proof}
Cut the cyclic order immediately after $a_0$.  For $r=0$ the
comparisons in~\eqref{gen:order-chain} say
$a_1<a_2<\cdots<a_{k-1}$; for $r=1$ they say the reverse chain.
Either chain determines the complete order of the remaining labels.
\end{proof}

For every occurrence of $\mu$ add the $k-2$ rows
\begin{equation}\label{gen:chain-rows}
 \orcyc(a_0,a_i,a_{i+1})
   =\orcyc_{\pi_\mu}(a_0,a_i,a_{i+1})\oplus r_\mu
 \qquad (1\le i\le k-2).
\end{equation}
At the median of any such triple in the represented embedding tree,
there are either three rigid R-directions or three marked P-directions.
The former gives a native affine flip expression; the latter gives its
complete marked frame plus a reference constant.  A P median for three
common pole-to-pole directions has the same other pole $q$: deleting
any other proposed pole would leave at least two of the three directions
connected through $q$.  In particular, this read is never a vacuous
two-anchor observation.

The chain also synchronizes the frame of every colour P-source carrying
at least three common direction groups.  Indeed, select the group
containing $a_0$.  Among the remaining groups some boundary
$a_i,a_{i+1}$ separates two successive groups.  The corresponding
chain triple uses three distinct source directions, so its equation
directly relates the source frame to $r_\mu$.  Thus all relevant source
frames are correlated with the common order, not merely its projection
on one chosen triple.

By Lemma~\ref{gen:chain-lemma}, the rows force exactly the prescribed
reference order or its reversal.  Together with the common R rows they
synchronize the entire oriented embedding of $K$ through the SPQR
parameterization.  S-skeletons have no independent choice.  Copy
reference signs and pole conventions are handled by~\eqref{gen:copy-sign}.

\subsection{Localization of the only large-domain side read}

Fix a fundamental cycle basis $\mathcal C(K)$.  For a connected set $Y$
outside a basis cycle $C$ in a represented block $B$, let $W$ be the
component of $B-V(C)$ containing $Y$.  Select an edge $ax$ with
$a\in C$ and $x\in W$.  The local gap of $ax$ between the outgoing and
incoming darts of $C$ is exactly $\side(C,Y)$, for every such edge,
because $W$ is connected and avoids $C$.  Any choice may therefore be
fixed, for example the first such edge in a fixed order; the statements
below hold for every choice.  Read this triple at its embedding-tree
median.

An R median gives a native flip plus a constant.  A two-marked-direction
P median gives a two-sector category, its frame, and a constant.  If a
P median has at least three marked directions and the observed third
direction is marked, the triple is determined by its complete marked
frame.  Only an unmarked third direction at such a P-source would produce
a noninjective interval predicate on an unrestricted larger category.

\begin{lemma}[Exceptional-read localization]
\label{gen:localization}
In that last case, the source poles $p,q$ are the poles of a common
P-node, the observed expansion $J$ is disjoint from $K$ except at its
poles, and
\[
       Y=W=J-\{p,q\},\qquad A_i(K,Y)=\{p,q\}.
\]
\end{lemma}

\begin{proof}
The cycle $C$ uses two distinct source expansions, so it traverses both
poles.  Lemma~\ref{gen:groups} gives the common P-node.  Every expansion
meeting $K$ outside the poles contains a whole common bridge
(Lemma~\ref{gen:groups}), hence a $K$-dart at the median's vertex.  The
marked subtree of Section~\ref{gen:marking} is the minimal subtree spanning
all $K$-dart leaves, so every such expansion is marked.  Therefore the
unmarked expansion $J$ contains no vertex of $K$ except its poles.

Its interior is connected and cannot leave to another expansion after
the deletion of $C$, because both poles were deleted.  Hence the
component reached by the selected attachment edge is exactly
$W=J-\{p,q\}$.  As $W$ contains no common-block vertex, it is also
the entire component $Y$ of $B-V(K)$ that it meets.  Both poles are its
attachments.
\end{proof}

This statement identifies the actual original lobe, not a coarse
projection of some other component.  That distinction is what permits
the domain restriction in the next step.

\subsection{Face witnesses and fixed basis signatures}

\begin{lemma}[Auxiliary hub witnesses]
\label{gen:hub}
Let $A\subseteq V(K)$ contain at least two distinct vertices, and let
$X$ be $K$ with a new vertex $z$ adjacent exactly to $A$.  Then $X$ is
biconnected.  For any fixed oriented embedding $\kappa$ of $K$, it has
an embedding extending $\kappa$ if and only if $A$ is incident with one
face of $\kappa$.  Every such face can contain the new vertex in an
extension, and in every extension the hub face is incident with all of
$A$.
\end{lemma}

\begin{proof}
Deleting $z$ leaves $K$ connected.  Deleting a vertex of $K$ leaves
$K$ minus that vertex connected, and $z$ retains a neighbour because
$|A|\geq2$.  This proves biconnectivity.  In an extension, the hub and
the interiors of its incident edges form a connected set disjoint from
$K$, so they lie in one face and every neighbour is on that face's
boundary.  Conversely, the boundary of a face of the biconnected plane
graph $K$ is a simple cycle.  Put $z$ inside its disk and draw a star
to the specified boundary vertices in their boundary order.  Its edges
are disjoint except at $z$, giving the required extension.
\end{proof}

\begin{lemma}[Fixed basis signatures distinguish faces]
\label{gen:basis}
Let $\mathcal C(K)$ be any fixed graph cycle basis consisting of simple
cycles.  In every planar embedding of the biconnected graph $K$, distinct
faces have distinct side vectors with respect to $\mathcal C(K)$.
\end{lemma}

\begin{proof}
Choose a curve between points in two distinct faces, transverse to edges
and avoiding vertices.  For a simple graph cycle, the parity of its
crossings with the curve is the XOR of the endpoint side bits.  This
crossing parity is linear on the binary edge-incidence vectors of cycles.
If both faces had the same basis signature, it would vanish on the whole
cycle space.  But the boundary of the first face is a simple cycle: the
first face lies on one side and every other face lies on the other.
That boundary cycle has odd crossing parity, a contradiction.  The proof
uses the same graph basis in every embedding and does not assume that
its cycles remain facial as source choices change.
\end{proof}

\subsection{Pairs with an extra attachment: a binary restriction}

Call $(K,Z)$ a \emph{full-sector pair} if $A(K,Z)$ is exactly the two
poles of a common P-node.  For every other pair, retain an independent
auxiliary graph
\[
 X(K,Z)=K+z_{K,Z},\qquad N(z_{K,Z})=A(K,Z).
\]
Test its planarity before constructing source references.  By
Lemma~\ref{gen:hub}, it is biconnected, and it extends a specified common
embedding exactly when all attachments share one face.  Every eligible
face can contain its hub.  Its comparison copy $K'$ is synchronized by
the common source rows above.

\begin{lemma}[Two relevant categories]
\label{gen:two-categories}
For every exceptional read of a lobe $Y\subseteq Z$ in a pair that is
not a full-sector pair, there is a fixed computable unary list of at
most two complete source categories which contains its category in
every \SEFE.
\end{lemma}

\begin{proof}
Lemma~\ref{gen:localization} supplies a common P-node $\mu$ with
$A_i(K,Y)=\{p,q\}\subseteq A(K,Z)$.  Because the pair is not a
full-sector pair, there is an extra attachment
$a\in A(K,Z)-\{p,q\}$.  Select any one; every choice gives a valid list,
and for definiteness the construction uses the first in a fixed vertex
order.  It belongs to the interior of a unique
common $\mu$-direction $e$.  A face incident with $p,q,a$ must be
one of the two reference sectors adjacent to $e$.  These names are
fixed even when internal reflections change their boundary paths or
whether $a$ is incident with them.

The exceptional source has at least three marked directions.  Its
source-to-common-sector map is the injection of
Lemma~\ref{gen:complete-read}.  Pull back the two sectors adjacent to
$e$ through that injection.  The resulting list has size at most two.
\end{proof}

Add the list as a unary restriction.  An empty list rejects; a singleton
pins the category.  For a two-element list $(c_0,c_1)$ introduce a Boolean
$b$ and the partial bijection $b\mapsto c_b$.  Every interval predicate
on that restricted category is a constant, $b$, or $b\oplus1$.
Consequently the original gap read becomes affine in $b$ and its source
frame.  This is a two-entry truth-table substitution, not an affine
representation of arbitrary triple orientations.

The hub has no exceptional large-domain read.  (Its reads lie in the
comparison of $X(K,Z)$ with $K'$, which is marked as in
Section~\ref{gen:marking}.)  If its attachment set has
size at least three, Lemma~\ref{gen:localization} would contradict that
size.  If it has size two and that pair is not a common P-pair, the same
lemma contradicts the absence of that common P-node.

For every lobe $Y\subseteq Z$ and every $C\in\mathcal C(K)$ retain the
original equality
\begin{equation}
 \side(C,Y)=\side(C,z_{K,Z}).
 \label{gen:side-equality}
\end{equation}
Every term is now affine.  The unary list is used only as a necessary
condition in the forward direction.  In the inverse direction no face
condition is discarded: every equality~\eqref{gen:side-equality} remains,
and the substituted expression is exactly the original side bit on its
restricted domain.

\subsection{Full-sector pairs: direct equality of the whole face}

For a full-sector pair $(K,Z)$ at common P-node $\mu$, introduce a variable
$t_{K,Z}\in\{0,\ldots,k-1\}$.  No auxiliary hub or cycle-basis side row
for this pair is needed.  Every colour lobe in $Z$ has at least two
attachments, all in $\{p,q\}$, so its attachment set is exactly that
pair.  Use Lemma~\ref{gen:complete-read} to obtain its exact face variable
$t_Y$, and impose
\begin{equation}
                    t_Y=t_{K,Z}.
 \label{gen:whole-face}
\end{equation}
This is a bijection.  A constant read pins $t_{K,Z}$; a two-category read
uses the affine gauge equation~\eqref{gen:two-gauge} followed by its
injection.  If there are no colour lobes, the sector is free.  Every named
sector is incident with both attachment vertices, so the attachment-face
condition is automatic.

Equation~\eqref{gen:whole-face} retains the complete face choice.  In
particular, the two ends of a union factor cannot independently choose
different sectors merely because each endpoint choice is legal.
Appendix~\ref{ver:two-pole} gives an instance showing that this relation
cannot be dropped.

\subsection{The ordinary graph and the marked-block invariant}
\label{gen:marking}

Take all source-labelled embedding trees of all original colour blocks,
including vertices without common darts.  Correlate each R-source's
copies by a Q-cell and the two P-pole copies by one pipe.  Join matching
red and blue common-dart leaves at every original vertex, with the reversed
blue convention.  Lemma~\ref{rep:lem:two-partitions} makes planarity of this
ordinary doubled graph equivalent to condition (O).

For each auxiliary graph, make an independent doubled comparison with
its literal common copy $K'$.  Include the hub's source tree despite its
lack of common joining leaves.  Source synchronization ensures that this
copy has the common embedding chosen by the originals.

For each nontrivial common block $K$ at a vertex $v$, mark the
minimal subtree spanning all $K$-dart leaves in each of the two
represented embedding trees, together with the matching leaf joins.
A marked piece is therefore two trees connected at the same leaf set,
which contains at least two darts.  It is biconnected: after deleting
an internal node of one tree, each remaining branch contains a marked
leaf and the other tree connects those leaves; deleting a joined leaf
leaves connection through the other leaves.  Common bridges remain
unmarked.  This construction is independent of the later choice of a
cycle basis for the face comparisons.

Mark every auxiliary comparison in the same way.  In the comparison of
$X(K,Z)$ with $K'$, at each $v\in V(K)$ mark the minimal subtrees spanning
the $K$-dart leaves in the trees of $X(K,Z)$ and of $K'$, together with
their leaf joins.  Only $K$ is marked there, and the hub's darts are
unmarked.  The groups at a P-source of $X(K,Z)$ are intervals of $\pi_\mu$,
or the instance is rejected (Lemma~\ref{gen:groups} applies to $X(K,Z)$),
and the marked P-copies of the comparison receive reference lists
projected from $\pi_\mu$ as below.  The following invariant argument
applies verbatim to these comparisons.

Every marked edge belongs to a cycle.  A free P-copy can meet only one
such common-block piece, by the two-poles argument of
Lemma~\ref{gen:groups}.  Its marked edges therefore lie in one
biconnected marked block.  A source pipe maps exactly the marked source
directions at both poles, since the marked $K$-directions agree at the two poles.  Thus pipes are off at both ends or aligned
on their whole marked sets.  These are precisely the marking hypotheses
of the complete-sector engine.

At a marked P-copy with at least three marked directions, use the
reference order of Lemma~\ref{gen:groups}.  At a two-marked-direction copy
choose an ordered pair $(a_0,a_1)$ and keep its frame explicit; the
constant $\delta$ in~\eqref{gen:two-gauge} refers to this choice.  All observations above are at
actual source copies, so no additional pipe membership or transparent
replication gadget is assumed.

\subsection{Correctness and witness transformations}

\begin{proof}[Proof of Theorem~\ref{gen:compilation}]
Suppose first that a \SEFE{} exists.  Restrict it to its colour blocks and
choose their exact SPQR source rotations.  These satisfy the ordinary
graph, pipes, Q-cells, and (O).  Set the common R and P variables from the
common embedding.  All native-R and common-P triple rows hold.

For a pair that is not a full-sector pair, the connected drawing of $Z$
lies in one face of $K$ incident with every attachment.  Embed its
independent auxiliary star in that face and give $K'$ the same common
embedding.  Every side equality holds.  The unary two-category lists are
satisfied by Lemma~\ref{gen:two-categories}; choose each Boolean selector
to record the observed category.  Its affine read is then the original
side bit.

For a full-sector pair, set $t_{K,Z}$ to the actual named sector of its
face.  Every lobe occupies that same face, and
Lemma~\ref{gen:complete-read} gives its exact category or constant,
including the two-anchor gauge.  Thus every whole-face equality holds.
This extends every solution; it makes no canonical choice of a private
order or an interleaving at a common cutvertex.

Conversely, decode a target witness into planar embeddings of all
original colour blocks and auxiliary graphs.  The original doubled graph
gives (O).  The common R rows and P-order chain rows identify the whole
oriented embedding of each common block in both colours and every
auxiliary copy.

For a pair that is not a full-sector pair, choose the hub's face as
$F(K,Z)$.  It is incident with all attachments.  Every substituted affine
read is the actual original side bit, including the two-entry exceptional
truth tables.  Equation~\eqref{gen:side-equality} therefore gives each
lobe the hub's full cycle-basis side vector.  By
Lemma~\ref{gen:basis}, such vectors distinguish faces in every embedding.
Hence every lobe lies in
$F(K,Z)$.

For a full-sector pair, interpret $t_{K,Z}$ as its unique pole-to-pole
face using Lemma~\ref{gen:p-faces}.  It is incident with both attachments.
The exact complete-face read and equation~\eqref{gen:whole-face} put
every lobe in that face.

Thus all selected faces satisfy every attachment incidence and every
lobe pin of Theorem~\ref{rep:thm:repair}.  Apply that theorem to the decoded
block embeddings and (O).  It produces a \SEFE, with all relative
component positions, common cutvertices, bridges, and isolated common
vertices accounted for.  The normalization inverse reconstructs a
solution of the original instance.  All operations are explicit source
decoding, face traversal, or the polynomial repair construction.  The
size and time bounds are derived in Section~\ref{gen:size}.
\end{proof}

\subsection{Size of the compiled instance}
\label{gen:size}

Let $N$ be the size of $I$.  We count the explicit data.
\begin{enumerate}
\item \emph{Pairs and lobes.}  By the charging argument of
      Section~\ref{rep:algorithm}, there are $O(N)$ pairs $(K,Z)$ and
      $O(N)$ colour lobes in total.  Also $\sum_K|E(K)|\le N$.
\item \emph{Graphs.}  The original colour blocks have total size $O(N)$.
      There are $O(N)$ auxiliary graphs $X(K,Z)$, one per pair, and each of
      them and its copy $K'$ has size $O(N)$, so all auxiliary graphs have
      total size $O(N^2)$.  An embedding tree at
      a vertex has one leaf per dart, and its inner nodes have degree at
      least three, so the source-labelled embedding trees of a graph have
      size linear in that graph.  Hence the ordinary doubled graph, its
      pipes and Q-cells, the marking, and all reference lists have total
      size $O(N^2)$.
\item \emph{Rows.}  A block $K$ has $|\mathcal C(K)|\le|E(K)|$ basis
      cycles and $O(N)$ pairs and lobes, so the side
      equalities~\eqref{gen:side-equality} number
      $O(N\sum_K|E(K)|)=O(N^2)$.  The block $K$ occurs in $B_1(K)$,
      $B_2(K)$, and in $X(K,Z)$ and $K'$ for each of its $O(N)$ pairs; each
      occurrence receives $O(|E(K)|)$ chain and R rows, so there are
      $O(N^2)$ such rows.  The equalities~\eqref{gen:whole-face} number
      $O(N)$.  Each row has constant width, with at most one Boolean
      selector and one unary list of size at most two per exceptional
      read.
\item \emph{Finite domains.}  Every category observer and every variable
      $t_Y,t_{K,Z}$ has a domain of size at most $N$, and every partial
      bijection has an explicit table of at most $N$ pairs.  There are
      $O(N^2)$ observers and relations, so domains and tables have total
      size $O(N^3)$.
\end{enumerate}
Altogether the output has size $O((N+1)^3)$.  Each item is produced by a
search in one graph or embedding tree of size $O(N)$, after planarity
tests and SPQR decompositions of the $O(N)$ graphs involved.  Even with a
fresh $O((N+1)^2)$-time computation for each of the $O(N^2)$ items, the
construction takes $O((N+1)^4)$ time.

